\section*{第 \ref{ch:g8:balanced-equations} 章——质量守恒与配平的化学方程式}

\begin{solution}{exo:g8:balanced-equations:1}
\omterm{def:g7:chemical-reactions:reaction}{化学反应}中，生成物的总质量
等于参加反应的\omterm{def:g7:chemical-reactions:reactant}{反应物}的总质量。
\end{solution}

\begin{solution}{exo:g8:balanced-equations:2}
\omterm{def:g8:balanced-equations:coefficient}{化学计量数}为 3。它表示 3 个碳\omterm{def:g7:atoms-and-molecules:atom}{原子}和 $3 \times 2 = 6$
个氧\omterm{def:g7:atoms-and-molecules:atom}{原子}。
\end{solution}

\begin{solution}{exo:g8:balanced-equations:3}
是的：左右两边各有 1 个 C 和 2 个 O。
\end{solution}

\begin{solution}{exo:g8:balanced-equations:4}
\ce{2Mg + O2 -> 2MgO}：左右两边各有 2 个 Mg 和 2 个 O。
\end{solution}

\begin{solution}{exo:g8:balanced-equations:5}
氮：左边有 2 个 N，因此取 2 \ce{NH3}；右边有 6 个 H，因此取 3
\ce{H2}：\ce{N2 + 3H2 -> 2NH3}。
\end{solution}

\begin{solution}{exo:g8:balanced-equations:6}
由\omterm{def:g8:balanced-equations:conservation}{质量守恒}定律可得：$44 - 12 = \qty{32}{g}$ 的氧气。
\end{solution}

\begin{solution}{exo:g8:balanced-equations:7}
生成的气体逸入房间，天平不再称量它。
没有质量消失：缺少的质量在房间的\omterm{def:g4:air-a-mixture-of-gases:air}{空气}中。
\end{solution}

\begin{solution}{exo:g8:balanced-equations:8}
碳：左右各有 1 个。氢：左边有 4 个，因此取 2 \ce{H2O}。氧：
右边有 $2 + 2 = 4$ 个，因此取 2 \ce{O2}：
\ce{CH4 + 2O2 -> CO2 + 2H2O}。
\end{solution}

\begin{solution}{exo:g8:balanced-equations:9}
把 \ce{H2O} 改为 \ce{H2O2} 改变了物质：\ce{H2O2} 是
过氧化氢，而不是水。只能改变\omterm{def:g8:balanced-equations:coefficient}{化学计量数}：
\ce{2H2 + O2 -> 2H2O}。
\end{solution}

\begin{solution}{exo:g8:balanced-equations:10}
$20 \times 5 = 100$ 个氧气\omterm{def:g7:atoms-and-molecules:molecule}{分子}；$20 \times 4 = 80$ 个水
\omterm{def:g7:atoms-and-molecules:molecule}{分子}。
\end{solution}

\begin{solution}{exo:g8:balanced-equations:11}
取 2 \ce{C2H6}：左边有 4 个 C，因此取 4 \ce{CO2}；有 12 个 H，因此取 6 \ce{H2O}；右边
氧\omterm{def:g7:atoms-and-molecules:atom}{原子}数为 $8 + 6 = 14$，因此取 7 \ce{O2}：
\ce{2C2H6 + 7O2 -> 4CO2 + 6H2O}。
\end{solution}

\begin{solution}{exo:g8:balanced-equations:12}
仍为 \qty{850.0}{g}。盒子密封：\omterm{def:g4:what-a-fire-needs:burning}{燃烧}掉的 \qty{0.4}{g} 蜡变成了二氧化碳和水蒸气，
它们与所吸收的氧气一起仍留在盒内。
没有物质进入或离开。
\end{solution}

\begin{solution}{pb:g8:balanced-equations:1}
\textbf{1.} 来自\omterm{def:g4:air-a-mixture-of-gases:air}{空气}中的氧气，它与铁结合：
氧化铁的质量等于铁加上这些氧气的质量。

\textbf{2.} 所有物质都留在广口瓶内：所用的氧气原本就在瓶中，
生成的氧化物也留在瓶中，总质量
不变。

\textbf{3.} \omterm{def:g8:balanced-equations:conservation}{质量守恒}：\omterm{def:g7:chemical-reactions:reaction}{化学反应}中，生成物的总质量
等于\omterm{def:g7:chemical-reactions:reactant}{反应物}的总质量。

\textbf{4.} \ce{4Fe + 3O2 -> 2Fe2O3}。

\textbf{5.} \ce{3Fe + 2O2 -> Fe3O4}。

\textbf{6.} 左边：4 个 Fe，$3 \times 2 = 6$ 个 O。右边：$2 \times 2 = 4$ 个 Fe，
$2 \times 3 = 6$ 个 O。平衡。

\textbf{7.} 4 个铁\omterm{def:g7:atoms-and-molecules:atom}{原子}对应 3 个氧\omterm{def:g7:atoms-and-molecules:molecule}{分子}，因此 $400 \times
\frac{3}{4} = 300$ 个氧\omterm{def:g7:atoms-and-molecules:molecule}{分子}。

\textbf{8.} $0.28 \times \frac{3}{7} = \qty{0.12}{g}$ 的氧气。

\textbf{9.} $0.28 + 0.12 = \qty{0.40}{g}$ 的氧化铁。

\textbf{10.} 是的：需要 \qty{0.12}{g}，而广口瓶中约有
\qty{0.5}{g}。

\textbf{11.} 所用的氧气已离开广口瓶内的\omterm{def:g4:air-a-mixture-of-gases:air}{空气}（现存于固态氧化物中），
因此瓶内气体比原来少，房间里的\omterm{def:g4:air-a-mixture-of-gases:air}{空气}流入。
天平读数随之上升，上升量等于进入的
\omterm{def:g4:air-a-mixture-of-gases:air}{空气}的质量。

\textbf{12.} \qty{0.12}{g} 的氧气取自广口瓶的\omterm{def:g4:air-a-mixture-of-gases:air}{空气}中。
\end{solution}
